\section{Introduction}
\label{sec:introduction}

Complex coupled systems---data centers, microgrids, process plants,
distributed databases, microservice architectures, compiler
pipelines---share a fundamental problem: their internal coupling
structure is hidden. Which subsystems affect which others,
through which channels, at what strength? The design documentation
specifies intended interfaces, but operational reality includes
emergent coupling paths created by the shared substrate that no
blueprint captures and no expert enumerates~\citep{ljung1999sid}.

Three existing approaches all fail to discover this structure in live
systems. Analytical modeling requires knowing the system's structure
a priori. Statistical observation---cross-correlation, Granger
causality, transfer entropy---cannot separate causal coupling from
co-movement when both subsystems are actively adjusting, because the
dominant noise source is the agents' own
exploration~\citep{sugihara2012ccm, runge2018confounding}. Expert
systems capture design intent, not operational reality, and go stale
as the system drifts.

\paragraph{The confounding problem.}
Consider multiple autonomous agents, each independently tending a
portion of a shared substrate. Each agent adjusts its own parameters
toward its own objective. The agents are unaware of each other.
Yet they communicate involuntarily through the substrate: when one
agent changes its parameters, the system responds in ways
that shift other agents' readings.

When all agents are actively tending, their objective readings move
for two confounded reasons: their own parameter changes and the other
agents' influence. No passive observer can separate these
contributions. The confounding is structural, not statistical---no
amount of observational data resolves it without intervention. This
is why cross-correlation, Granger causality, and transfer entropy all
fail on actively explored systems. They cannot distinguish ``shifting
together because coupled'' from ``shifting together because each is
independently adjusting.''

\paragraph{The key insight.}
The confound disappears if one agent stops exploring and holds still
while another perturbs. With the receiver frozen, any change in its
signal is attributable to the sender's perturbation---not to its own
adjustments. The agent ceases to be an explorer and becomes a
sensor. The sender ceases to explore and becomes a stimulus.

This is the active probing paradigm, borrowed from radar and
seismology: inject a known signal into the medium and observe the
response. The system describes itself through its response to
perturbation, rather than through passive observation of its
behavior. The coupling graph is not inferred from correlations---it
is measured as a set of impulse responses.

\paragraph{Our approach.}
We embed this paradigm inside the agents' normal operation through a
coordinated push/pause/hold protocol. The protocol directs the agents
to take turns: one agent (the sender) drives its parameters to
extremes, then returns to neutral, while all others hold fixed as
passive sensors. An adaptive calibration phase finds a safe
perturbation amplitude before each push block, backing off
exponentially if the system destabilizes. A constant-false-alarm-rate
(CFAR) step-change detector evaluates whether each receiver's signal
exhibits a reversible shift---rising during push, recovering during
pause. The ABA block design eliminates drift confounds: a monotonic
trend produces a rising edge but not a falling edge.

The protocol requires no information exchange beyond scheduling
(``whose turn is it''), enforced through a shared lease with fencing
tokens. Agents do not share objectives, results, or detected edges.
Each discovers coupling purely by watching its own signal shift when
another agent perturbs the substrate.

\paragraph{Coupling-shape agnosticism.}
Because the protocol measures step response---not frequency response,
not correlation structure---it is agnostic to the coupling function's
shape. Linear coupling produces a proportional step. Threshold
coupling produces a discontinuous jump. Saturation coupling produces
a compressed or inverted step. Polynomial coupling produces a
quadratic step. The detector asks only: did the signal change and
recover? The shape of the response determines the step amplitude but
not whether coupling exists.

\paragraph{Contributions.}
\begin{enumerate}
  \item A coordination protocol---push, pause, hold, with adaptive
    perturbation calibration and turn-taking through a fencing-token
    lease---that enables autonomous agents to discover coupling
    structure through their shared substrate without exchanging
    information beyond scheduling.
  \item Empirical characterization on a synthetic bench with planted
    ground truth across four coupling shapes (linear, threshold,
    saturation, polynomial), five noise levels, and six coupling
    strengths. Detection floor, noise boundary, and false positive
    rate are measured and reported honestly.
  \item Step amplitude measurements within 5\% of ground truth across
    the detectable range, with zero false positives across all noise
    levels tested.
\end{enumerate}

The contribution is the protocol that makes coupling
\emph{detectable}---not a novel detector. CFAR is established
technology from radar signal processing~\citep{finn1968cfar}. What
is new is the choreography: autonomous agents, unaware of each other
and communicating only through the substrate they tend, taking turns
as active perturbers and passive observers. The agent that was
exploring blindly becomes a measurement instrument. The probe
\emph{is} the experiment.
